% S2WISH technical report, AIMS Predictive AI Evaluation Competition (NeurIPS 2026)
% Build:  build.ps1 (PowerShell). It produces two PDFs from this one source:
%   main.pdf            submission: main text, appendices A-D, NeurIPS checklist (no atlas)
%   main_reference.pdf  reference:  the same plus the visual atlas
% The reference build is a two-line wrapper that defines \buildreference and inputs this file.
% Source of every number: docs/publication/FACTS_2026-10-16.md and the raw files it names.
\documentclass{article}

% Non-anonymous preprint layout. Switch to [final] if the organisers ask for camera-ready style.
\PassOptionsToPackage{numbers,sort&compress}{natbib}
\usepackage[preprint]{neurips_2026}

\usepackage[utf8]{inputenc}
\usepackage[T1]{fontenc}
\usepackage{hyperref}
\usepackage{url}
\usepackage{booktabs}
\usepackage{array}
\usepackage{amsfonts,amsmath}
\usepackage{nicefrac}
\usepackage{microtype}
\usepackage{xcolor}
\hypersetup{colorlinks=true,linkcolor=blue!50!black,citecolor=green!35!black,urlcolor=blue!50!black}
\usepackage{float}
\raggedbottom
\usepackage{graphicx}
\usepackage{subcaption}
\usepackage{enumitem}
\usepackage{siunitx}
\sisetup{group-separator={,},group-minimum-digits=4}

% Build switch: \referencetrue adds the visual atlas (reference PDF); the submission PDF leaves it out.
\newif\ifreference
\ifdefined\buildreference \referencetrue \else \referencefalse \fi

\newcommand{\dalc}{\ensuremath{\Delta\mathrm{ALC}}}
\newcommand{\ee}[1]{\ensuremath{\times 10^{-#1}}}
\newcommand{\rs}{\ensuremath{\rho_S}}
\newcommand{\fillin}[1]{\textcolor{red}{[FILL IN: #1]}}
% "Sources:" line under each table and figure: generating script and results file (paths relative to the repository).
\newcommand{\src}[1]{{\footnotesize\ Sources: #1.}}
\newcommand{\p}[1]{\nolinkurl{#1}}
% Pointer to the feedback-inversion appendix (appendix_inversion.tex, written separately). It resolves to
% "Appendix X" when that file defines \label{app:inversion}, and to plain words otherwise, so it never prints "??".
\makeatletter
\newcommand{\invref}{\@ifundefined{r@app:inversion}{the feedback-inversion appendix}{Appendix~\ref{app:inversion}}}
\makeatother
\graphicspath{{figures/}}
% Figures come from make_figures.py; a missing one renders as a labelled box so the text can be checked first.
\newcommand{\fig}[2][\linewidth]{\IfFileExists{figures/#2}{\includegraphics[width=#1]{#2}}%
  {\fbox{\parbox[c][1.6in][c]{0.9#1}{\centering\small missing: #2}}}}

\title{Predicting AI Performance on Unseen Benchmarks:\\
Shrinkage, Hierarchical Refinements and Their Limits}

\author{Shashwat Srivastava\\ Accenture\\ \texttt{shashwat1322001@gmail.com}\\ Team S2WISH\\
  Code: \url{https://github.com/Shashwat-srivastav/s2wish-paiec2026}}

\begin{document}
\maketitle

\begin{abstract}
On benchmarks a model has never seen, shrinkage is the working core; what remains is item difficulty, which none of
the label-free signals we tested without a runtime LLM supplies usefully (runtime LLM priors, which the rules permit,
were not tested). We show this for the AIMS Predictive AI Evaluation Competition, which scores
predictions on a hidden benchmark by the Brier area under the learning curve (ALC; lower is better). In leave-one-benchmark-out evaluation, empirical-Bayes shrinkage of early estimates gives
33.7 of the 36.3 thousandths that our small hierarchical Bayesian model gains over the official baseline. A
pre-registered check on 89 unseen public benchmarks confirms the shrinkage gain; there the check detected no
difference between the hierarchical model and shrinkage alone (inconclusive at the pre-registered $\pm1\ee{3}$), so
the refinements are unproven. On
the platform draws we observed, no pair-level adjustment could improve our model by more than about 0.022 ALC (95\%
CI up to 0.033), in-sample. We contribute floor and ceiling bounds that locate the remaining
loss in item difficulty; error bounds and exact replays for reading the platform's per-pair feedback; and negative
results from 60 logged experiments (22 pre-registered).
\end{abstract}

% ---------------------------------------------------------------------------------------------------------------
\section{Introduction}
\textbf{On unseen benchmarks, shrinkage does almost all the work; what remains is item difficulty, which none of the
label-free signals we tested without a runtime LLM supplies usefully.} Runtime LLM priors, which the rules permit,
were not tested. Figure~\ref{fig:teaser} shows the story in one picture.

Predicting outcomes is cheaper than running every AI system on every benchmark. The AIMS Predictive AI Evaluation
Competition~\cite{aims2026,truong2026} poses this under the hardest condition: the benchmark is new, so no system's
responses to its items are known. Item response theory (IRT) methods need exactly those responses
(\S\ref{sec:related}). We therefore built a model that learns only what transfers across benchmarks.

\noindent\begin{minipage}{\linewidth}
\paragraph{What we show.}
\begin{itemize}[leftmargin=*,nosep]
  \item \textbf{Shrinkage does almost all the work} (\S\ref{sec:results}). The hierarchical refinements are not
    established across benchmarks; a pre-registered out-of-sample check detects no difference from shrinkage.
    Table~\ref{tab:tried} lists every rejected candidate.
  \item \textbf{The remaining loss is item difficulty} (\S\ref{sec:headroom}). Oracles and floor and ceiling bounds
    show that pair-level calibration cannot close the gap; no label-free signal we tested without a runtime LLM
    usefully supplies item difficulty.
  \item \textbf{Per-pair feedback can be read exactly} (\S\ref{sec:inversion}). We give error bounds, the reflection
    pitfall and exact replays (\invref); they reveal that platform pairs are harder than our local folds.
\end{itemize}
\end{minipage}

\begin{figure}[H]
  \centering
  \fig{fig_teaser.pdf}
  \caption{Local ALC with one subject per benchmark, mean over four held-out benchmarks. Beta-binomial shrinkage
  gives 33.7 of the $36.3\ee{3}$ gain over the official baseline; the hierarchical refinements add the rest. Below
  the dashed pair-constant floor $a(1-a)$, the score of knowing each pair's accuracy but nothing about its items,
  only item-level knowledge helps (\S\ref{sec:headroom}).
  \src{\p{make_figures.py}; \p{experiments/runs/review_round1/ladder_intervals.json} and \p{floor_ceiling.json}}}
  \label{fig:teaser}
\end{figure}

% ---------------------------------------------------------------------------------------------------------------
\section{Related work}\label{sec:related}
\textbf{IRT and efficient evaluation.} tinyBenchmarks~\cite{polo2024}, metabench~\cite{kipnis2025}, Anchor
Points~\cite{vivek2024}, Fluid Benchmarking~\cite{hofmann2025}, amortized calibration~\cite{truong2025} and IRT scaling
laws~\cite{truong2026irsl} all calibrate items on many systems' responses to the target benchmark, which a hidden
benchmark does not allow; contextual IRT with text embeddings ranks items at chance level across
scenarios~\cite{shang2026}, as our text features do (Table~\ref{tab:repr}). \textbf{Unseen data.} Instance-level
assessors degrade out of distribution~\cite{pacchiardi2024,pacchiardi2025}; forecasting benchmark accuracy from a task
description needs an LLM at prediction time~\cite{park2025}; system-level abilities do transfer across
benchmarks~\cite{ruan2024}. \textbf{Empirical Bayes.} Shrinkage of many small-sample rates toward a pooled level beats
the raw mean~\cite{efron1975,brown2008}, and heavy-tailed weakly informative priors are a robust default for logistic
models~\cite{gelman2008}. \textbf{Scoring and calibration.} The Brier score~\cite{brier1950} is strictly
proper~\cite{gneiting2007}; Murphy's decomposition~\cite{murphy1973} separates reliability from resolution; temperature
scaling~\cite{guo2017} is the standard recalibration. \textbf{Acquisition and noise.} Active testing~\cite{kossen2021}
and active inference~\cite{zrnic2024} reduce to uniform sampling when items cannot be told apart, and adaptive testing
on noisy item parameters can hurt~\cite{vanderlinden2000}; small-sample evaluations need wider intervals than the
central limit theorem suggests~\cite{bowyer2025,miller2024}, and repeated runs on the same item vary~\cite{wang2025}.

% ---------------------------------------------------------------------------------------------------------------
\section{Task, evaluation and data}\label{sec:task}
A \emph{subject} is one AI-system configuration, described by eight strings: model name, provider, release and access
dates, harness, harness version, reasoning effort and extras. An \emph{item} is described by four: content, features,
interactors and an anonymous benchmark ID. A submission implements \texttt{predict(input, labeled)}~$\to p$ and may
implement a label-acquisition hook. For each subject--benchmark pair the platform splits the items 50/50 into an
acquisition pool and an evaluation pool. It reveals up to 31 acquisition labels and scores the evaluation pool at six
label budgets:
\begin{equation}
\mathrm{ALC} = 0.1B_0 + 0.2B_1 + 0.2B_3 + 0.2B_7 + 0.2B_{15} + 0.1B_{31},
\end{equation}
where $B_n$ is the Brier score, the mean of $(p-y)^2$, with $n$ labels revealed.
ALC is averaged over pairs with equal weight, so always predicting 0.5 scores 0.25.

\newpage
At every budget,
\texttt{labeled} holds the revealed answers of \emph{every} sampled pair. Another subject's label on the same hidden
benchmark, or on the same item, is therefore legal evidence.

We re-implemented the evaluator and drove the organisers' unmodified streaming client through it. After we fixed nine
discrepancies in the first version, per-pair, per-budget Brier scores matched exactly (maximum difference 0). Each
formative submission scores a fresh draw of about \num{1000} items. In practice that is eight or nine pairs over seven
anonymous benchmarks, usually with one subject per benchmark.

\begin{table}[H]
  \caption{Public training data (revision \texttt{bc8204d811}). MMDocRAG (not binary) is unused; SWE-rebench (one
  subject) is used for training only.
  \src{\texttt{scripts/reproduce\_all.py -{}-only T1}, from the gated \p{data/*/response.parquet}}}
  \label{tab:data}
  \centering\small
  \begin{tabular}{lrrrrr}
    \toprule
    benchmark & subjects & items & responses & pairs $\ge$ 80 items & mean pair accuracy\\
    \midrule
    MathArena          & 104 & \num{1755} & \num{92944} & 81 & 0.61\\
    Multi-SWE-bench    &  82 & \num{2126} & \num{57975} & 82 & 0.22\\
    REAL web agents    &  33 &  233 &  \num{4449} & 26 & 0.34\\
    ResearchCodeBench  &  31 &  212 &  \num{6784} & 31 & 0.35\\
    SWE-rebench        &   1 & \num{6306} & \num{67074} &  1 & 0.48\\
    \bottomrule
  \end{tabular}
\end{table}

\paragraph{Validation protocol and notation.} Benchmarks are held out whole, so the only meaningful validation is
leave-one-benchmark-out (LOBO). That gives four folds (Table~\ref{tab:data}: Multi-SWE-bench, REAL, ResearchCodeBench,
MathArena), so every benchmark-level statement in this report rests on four numbers. Within a fold we draw test
episodes in one of three ways:
\begin{itemize}[leftmargin=*,nosep]
  \item \textbf{one subject per benchmark}, which is what the platform's draws turned out to be. This is our primary
    setting, and every local number is in it unless stated otherwise;
  \item \textbf{platform-like draw}: about \num{1000} items sampled as the platform does;
  \item \textbf{all subjects}: every subject present. This is an upper bound only, because it overstates
    cross-subject pooling about sevenfold.
\end{itemize}
All comparisons are paired on identical episodes. We quote differences in ALC in thousandths, written
\dalc{}\,\ee{3}: a value of $-1.0$ means 0.001 lower (better) ALC. Unless stated otherwise, intervals are 95\%
bootstrap intervals over pairs. \rs{} denotes Spearman rank correlation. Research notes cited by number are in
\texttt{docs/research\_notes/} of the code release.

\paragraph{Where the predictable signal is.}
Per-benchmark Rasch fits place more logit-scale variance on items than on subjects for every benchmark. Re-running the
same subject on the same item accounts for a further 20--27\% of response variance (repeated cells in MathArena and
SWE-rebench; note 11), which no predictor can remove. Subject metadata explains 16\% of ability variance on an unseen
benchmark (ridge $R^2$ under LOBO; note 14), mostly through release date. Model ranks transfer between coding
benchmarks (\rs{} 0.6--0.9), but levels do not (intraclass correlation 0.22). Text features rank items within a
benchmark but do not transfer to new benchmarks; an LLM difficulty rating, computed offline, transfers weakly
(Table~\ref{tab:repr}, Appendix~\ref{app:sources}). A leakage audit found no runtime-visible field that predicts outcomes beyond the benchmark
mean. It also showed that folds must be grouped by item content rather than item ID, because duplicates inflate
MathArena by 0.014 Brier.


% ---------------------------------------------------------------------------------------------------------------
\section{Method}\label{sec:method}
The predictor is a Rasch-type model with four latent quantities per prediction:
\begin{equation}
\operatorname{logit} P(y=1) = \mu_b + \theta_{sb} - \beta_i + \gamma_{g(i)},
\end{equation}
\vspace{-1.2em}
\begin{align*}
\mu_b &\sim t_3(G/2,\ s) && \text{benchmark offset; evidence: every subject's labels on benchmark } b\\
\theta_{sb} &\sim \mathcal{N}(0,\tau^2) && \text{pair ability; evidence: the subject's own labels}\\
\beta_i &\sim \mathcal{N}(0,\sigma_\beta^2) && \text{item difficulty; evidence: other subjects' labels on the same item}\\
\gamma_g &\sim \mathcal{N}(0,\lambda\sigma_\beta^2) && \text{item-group effect; evidence: the subject's labels in the same group.}
\end{align*}

Fitted on the five public benchmarks, $G=-0.58$ (the mean logit accuracy), $s=0.78$, $\tau=1.13$ and
$\sigma_\beta=2.70$; $\lambda$ is the share of item-difficulty variance assigned to groups. Four design choices
matter.
\begin{itemize}[leftmargin=*,nosep]
  \item \textbf{A hedged offset prior.} The offset prior is weakly informative and heavy-tailed~\cite{gelman2008},
    centred at $G/2=-0.29$, halfway between the pooled level and 50\%. This cuts worst-case regret on unusually easy
    benchmarks by 29\% (platform-like draws) to 50\% (one subject per benchmark) at no mean cost.
  \item \textbf{Numerical integration.} The posterior is integrated by deterministic numerical quadrature: a
    121-point grid over $\mu$ and Gauss--Hermite quadrature over $\theta$ and $\beta$, with no sampling. In an early development run on SWE-rebench, plug-in point
    estimates reached 0.49 Brier after one label (research record, overview of design decisions).
  \item \textbf{Item pooling.} Items that share a key are pooled across subjects through the false-merge model and
    the optional key guard (Table~\ref{tab:ladder}).
  \item \textbf{Item groups.} Groups come from the fields of \texttt{item\_features}, with $\lambda$ averaged over
    three values by marginal likelihood.
\end{itemize}
Predictions are clipped to $[0.005, 0.995]$. Labels are acquired at random. Four non-random policies gained at most
$0.27\ee{3}$, as theory predicts when the predictor cannot tell unlabelled items apart~\cite{kossen2021,vanderlinden2000}.
The archive is pure Python and NumPy (three files, about 5\,KB zipped) and takes 3.3--222\,ms per call.

We built the model as a ladder of versions (Table~\ref{tab:ladder}). The \emph{final model}, our selected submission, is the
item-group model with the key guard switched off; its two code files are byte-identical to the item-group model's.
Archive names and checksums are in Appendix~\ref{app:rerun}.

% ---------------------------------------------------------------------------------------------------------------
\section{Results}\label{sec:results}
\subsection{Local leave-one-benchmark-out}
\begin{table}[H]
  \caption{Brier score by label budget and ALC with one subject per benchmark, mean over the four held-out benchmarks
  (10 seeds, all rows paired; budgets 3 and 15 in Appendix Table~\ref{tab:ladder-full}). The official baseline is
  the empirical mean. Each later row adds: shrinkage; a benchmark offset and item pooling (first model); the hedged
  centre $G/2$ and the false-merge model (hedged-centre model); item groups and the key guard (item-group model); key
  guard off (final model). \emph{False-merge model}: items sharing a key are pooled, with a probability $r$ that
  the key wrongly merges items. \emph{Key guard}: sets $r\ge0.5$ when all of a key's labels are the target subject's.
  \emph{Step}: \dalc{}\,\ee{3} against the row above, with 95\% intervals from a hierarchical bootstrap (benchmarks,
  then pairs) and a $t$-interval over the four fold means; the last upper endpoint 0.0 is degenerate (two folds differ
  by exactly 0). The final-model row replays the shipped archive (platform-like draw: 0.2143).
  \src{\p{src/paiec/analysis/review_round1.py}; \p{experiments/runs/review_round1/ladder_intervals.json} and
  \p{table3_baseline_seeds0-9.json}}}
  \label{tab:ladder}
  \centering\small\setlength{\tabcolsep}{2.6pt}
  \begin{tabular}{@{}>{\raggedright\arraybackslash}p{0.27\linewidth}ccccc|rcc@{}}
    \toprule
    model & $B_0$ & $B_1$ & $B_7$ & $B_{31}$ & ALC & step & hierarchical & $t$, 4 folds\\
    \midrule
    official baseline & 0.250 & 0.360 & 0.209 & 0.189 & 0.2451 & & &\\
    Beta-binomial shrinkage & 0.250 & 0.231 & 0.201 & 0.188 & 0.2114 & $-33.70$ & $[-42.8, -23.4]$ & $[-50.6, -16.8]$\\
    first model & 0.247 & 0.229 & 0.200 & 0.187 & 0.2098 & $-1.62$ & $[-6.4, +6.4]$ & $[-14.1, +10.8]$\\
    hedged-centre model & 0.246 & 0.229 & 0.200 & 0.187 & 0.2097 & $-0.09$ & $[-5.2, +3.3]$ & $[-8.1, +8.0]$\\
    item-group model & 0.246 & 0.229 & 0.200 & 0.184 & 0.2089 & $-0.83$ & $[-1.4, -0.2]$ & $[-1.9, +0.2]$\\
    \textbf{final model} & 0.246 & 0.229 & 0.199 & 0.183 & \textbf{0.2088} & $-0.09$ & $[-0.3, 0.0]$ & $[-0.3, +0.1]$\\
    \bottomrule
  \end{tabular}
\end{table}
\textbf{Takeaway: shrinkage carries the gain; the refinements are small and not established at benchmark level.}
\emph{In-sample (four held-out benchmarks).} Shrinkage alone accounts for 33.7 of the total $36.3\ee{3}$ gain
(Table~\ref{tab:ladder}), as empirical Bayes predicts for small per-pair samples~\cite{efron1975,brown2008}. It is the
only step established at benchmark level. Everything added afterwards is worth $2.6\ee{3}$, concentrated at budgets
7--31, where item-level evidence accumulates. Its hierarchical interval, $[-4.5, -0.3]\ee{3}$, excludes zero; its
$t$-interval over the four fold means, $[-6.3, +1.1]\ee{3}$, includes zero. On MathArena, Beta-binomial shrinkage beats the
final model by $0.67\ee{3}$. Of the single steps, only the item-group step is supported, and only by the hierarchical
interval.

\emph{Out-of-sample (89 public benchmarks).} A pre-registered out-of-sample check uses 89 public benchmarks never used
for fitting (2{,}610 pairs, one subject per benchmark, frozen parameters, no item features). The same benchmarks were
used earlier to check the simulator and to reject the $\tau$/mixture variants, so earlier final-model scores on them
were known; the pre-registration discloses this. The check confirms the shrinkage step: $-32.7\ee{3}$ $[-37.1, -28.2]$ (benchmark
bootstrap), better on 82 of 89 benchmarks. Its primary outcome, the final model minus shrinkage, is $+0.2\ee{3}$
$[-0.9, +1.5]$: no difference detected, and inconclusive at the pre-registered $\pm1\ee{3}$ minimum effect.

We ship the hierarchical model because it was the pre-declared final model, chosen before the check ran; every paired
local step lowers the point estimate; and out of sample no difference from shrinkage was detected. Secondary and
descriptive results agree in direction but are not confirmatory. The platform-like (secondary) difference is
$-1.7\ee{3}$ $[-3.1, -0.2]$; of the 89 benchmarks, four derive from training benchmarks, and excluding them the upper
end reaches $+0.05\ee{3}$. Secondary outcomes never override the primary. Of the 89, the 5 with accuracy
$\le 0.15$ (a band holding about a third of the platform's pair-slots; Figure~\ref{fig:excess}) descriptively favour
the final model over shrinkage ($-7.1\ee{3}$; $-13.6\ee{3}$ platform-like); there, shrinkage alone is worse than the
official baseline by $+9.9\ee{3}$ $[-1.9, +22.8]$, on 3 of the 5.

\subsection{Official platform}
\begin{table}[H]
  \caption{Local versus platform. (a) Only the first final-model run has per-budget feedback; its SE is
  \emph{within-draw}, across pairs, not run-to-run. Over four runs the ALC is 0.1839 (SD 0.0240, SE 0.0120,
  $[0.146, 0.222]$); 1 of 4 runs lies above the local value (sign test $p=0.625$). (b) Platform: runs and mean with
  95\% $t$-interval (fresh draws, unpaired; the item-group archive only, without the re-package in
  Appendix Figure~\ref{fig:runs}). Replay: in-sample on the replay set, same labels for all versions; exact-set values in
  parentheses. The replay feeds labels without item features, so it cannot see the item-group step, and its order
  differs from the local column. On the exact set, baseline minus final is
  $+0.0165$ $[-0.0032, +0.0352]$ and shrinkage minus final $-0.0009$ $[-0.0052, +0.0029]$.
  \src{\p{src/paiec/analysis/review_round2.py}; \p{experiments/runs/review_round2/per_budget_final.json} and
  \p{local_vs_platform.json}; \p{experiments/runs/review_round1/platform_baseline_replay.json}}}
  \label{tab:platform}
  \centering\footnotesize\setlength{\tabcolsep}{4pt}
  \begin{tabular}{@{}lccccccc@{}}
    \toprule
    \multicolumn{8}{@{}l}{\emph{(a) final model, Brier by budget}}\\
    \midrule
    & $B_0$ & $B_1$ & $B_3$ & $B_7$ & $B_{15}$ & $B_{31}$ & ALC\\
    \midrule
    first final-model run      & 0.2369 & 0.1903 & 0.1613 & 0.1445 & 0.1463 & 0.1412 & 0.1663\\
    \quad within-draw SE (9 pairs) & 0.0106 & 0.0241 & 0.0268 & 0.0277 & 0.0292 & 0.0290 & 0.0243\\
    local, one subject per benchmark & 0.2456 & 0.2285 & 0.2123 & 0.1995 & 0.1894 & 0.1832 & 0.2088\\
    \bottomrule
  \end{tabular}\\[4pt]
  \begin{tabular}{@{}lccc@{}}
    \toprule
    \multicolumn{4}{@{}l}{\emph{(b) ALC per version}}\\
    \midrule
    version & local & platform: $n$, mean [95\% CI] & replay set (exact set)\\
    \midrule
    Beta-binomial shrinkage & 0.2114 & not submitted & 0.1978 (0.1888)\\
    first model             & 0.2098 & 1, 0.2157 & 0.1945 (0.1909)\\
    hedged-centre model     & 0.2097 & 5, 0.1923 $[0.176, 0.208]$ & 0.1954 (0.1896)\\
    item-group model        & 0.2089 & 4, 0.1909 $[0.170, 0.212]$\rlap{$^{\ast}$} & 0.1954 (0.1896)\\
    \textbf{final model}    & 0.2088 & 4, 0.1839 $[0.146, 0.222]$ & 0.1954 (0.1896)\\
    \bottomrule
  \end{tabular}\\[1pt]
  {\scriptsize $^{\ast}$Includes one run marked Failed but scored (0.1840; ID in Appendix~\ref{app:names}); without it,
  $n=3$, mean 0.1932.}
\end{table}
\textbf{Takeaway: on platform labels, a paired replay confirms the large step (baseline $\to$ shrinkage); the later
refinements are not distinguishable there, and their order differs from the local ladder.} Before its first
submission, the re-implemented evaluator predicted 0.212 for our first model, which scored 0.216. Same-archive SDs
are 0.010--0.024 across versions (0.013 for the hedged-centre model over 5 runs), of the size small-sample evaluation
theory predicts~\cite{bowyer2025,miller2024} and far larger than the differences between our versions
(Table~\ref{tab:platform}; every run in Appendix Figure~\ref{fig:runs}). The platform means do not resolve the order of the versions:
intervals overlap, and the order flips depending on whether the run marked Failed is counted. Model choices therefore
rest on the paired local tests.

\newpage
The in-sample replay (\S\ref{sec:inversion}) scores all versions on the same platform
labels of the replay set: the official baseline (0.2276) minus the final model is $+0.0322$ $[+0.0140, +0.0515]$,
while Beta-binomial shrinkage minus the final model is $+0.0023$ $[-0.0008, +0.0055]$. There, among the ladder versions, the first model (0.1945)
scores best, and on the exact set shrinkage (0.1888) beats the final model (0.1896).
Official levels sit below the local ones because platform pairs are more extreme \mbox{(median accuracy 0.27)}, which lowers
the irreducible $a(1-a)$ (\S\ref{sec:sim}). The final ranking uses one common hidden subset, so draw-to-draw
variation affects all submissions equally.

\subsection{What did not work}
\begin{table}[H]
  \caption{What we tried and why it was not shipped. Paired \dalc{}\,\ee{3} against the final model with one subject
  per benchmark, 95\% interval over pairs (positive is worse); the adoption gate is $\le -1.0$. ``Pre-reg.'': yes =
  pre-registered, post hoc = registered after the run. Appendix Figure~\ref{fig:forest} plots the first five rows.
  \src{\texttt{experiments/runs/preregistration/exp\_00\{42,46,52,53,57\}\_*.json}; \p{experiments/registry.csv}}}
  \label{tab:tried}
  \centering\footnotesize\setlength{\tabcolsep}{3.5pt}
  \resizebox{\linewidth}{!}{%
  \begin{tabular}{@{}llcl@{}}
    \toprule
    component & \dalc{} [95\% CI] & pre-reg. & decision and reason\\
    \midrule
    lexical residual smoothing & $+1.68$ $[+1.08, +2.30]$ & yes & dropped: worse\\
    gp-IRT, our reimplementation~\cite{polo2024} & $+6.88$ $[+4.80, +8.94]$ & yes & not promoted: worse\\
    re-tuned shrinkage strength & $+0.74$ $[-0.28, +1.85]$ & yes & no change: mild under-shrinkage only$^{a}$\\
    fitted temperature~\cite{guo2017} & $+0.40$ $[+0.24, +0.56]$ & yes & stopped: already calibrated$^{b}$\\
    symmetric mixture prior & $+0.52$ $[+0.22, +0.80]$ & yes & retired: worse\\
    item scorer & no paired ALC (failed its transfer gate) & yes & dropped\\
    label-free benchmark level & no paired ALC (failed its pre-registered stopping rule) & yes & dropped\\
    nested meta-learned priors & no paired ALC (failed its pre-registered gate) & yes & dropped\\
    prior factorial & no paired ALC (stopped after the pilot) & yes & Laplace-vs-grid audit open$^{c}$\\
    robust, metadata and model-name priors & best gain $\le 1.0$ & post hoc & below the gate\\
    smarter label acquisition (4 policies) & best gain $\le 0.27$ & post hoc & random acquisition kept\\
    LLM rating distilled into MiniLM/MPNet & no paired ALC (kept $\le$ 3\% of its \rs{}) & post hoc & dropped\\
    \bottomrule
  \end{tabular}}\\[1pt]
  {\scriptsize $^{a}$Empirical half-life 5.27 labels against 4.30 for the model. $^{b}$Fitted $T=0.945$--$1.05$.
  $^{c}$Differences 0.012--0.018.}
\end{table}
\textbf{Takeaway: nothing we tried beat the final model.} Table~\ref{tab:tried} lists every candidate with its paired
interval, gate and reason. Priors failed for one reason: one to three labels cannot tell prior centres a third of a
logit apart, and ability is mostly benchmark-specific. Acquisition failed because the predictor cannot tell
unlabelled items apart.

% ---------------------------------------------------------------------------------------------------------------
\section{Reading per-pair feedback exactly}\label{sec:inversion}
At budget 0 our model predicts one constant $p_0$ for every item of a pair. We obtain $p_0$ by calling each archive's
own \texttt{predict()} with no labels: 0.448722 for the hedged-centre, item-group and final models, and 0.400099 for
the first model. The pair's accuracy $a$ on its scored items then follows from the reported budget-0 Brier score:
\begin{equation}
B_0 = a(1-p_0)^2 + (1-a)p_0^2 \quad\Longrightarrow\quad a = \frac{B_0 - p_0^2}{1-2p_0}.
\end{equation}
The identity is elementary (the one-bin case of Murphy's decomposition~\cite{murphy1973}) and has been used
independently by others; we contribute error bounds, the reflection pitfall and exact replays (\invref).
Reported values are rounded to six decimals, with error at most $\eta=5\ee{7}$. The reported calibration error (ECE)
at budget 0 equals $|a-p_0|$. The two estimates of $a$ therefore agree to within $\varepsilon(p_0)=\eta/|1-2p_0|+\eta$;
on every saved pair they agree to within $5.1\ee{6}$. The floor $a(1-a)=B_0-(p_0-a)^2$ does not depend on choosing
between the two roots. The archive's $p_0$ is needed only to pick $a$ over $1-a$; getting it wrong flips every
accuracy (\S\ref{sec:sim}). On \emph{singleton} pairs, whose benchmark appears once in a draw, the prediction stays constant at every budget.
Brier score and ECE then fix the number of successes among the revealed labels, so a pair can be replayed exactly
in-sample. We use two sets: the \emph{replay set} (18 pairs that replay exactly over our five saved formative runs)
and the
\emph{exact set} (8 of them, where the item-group and final models also replay exactly; elsewhere item-group effects
break the count-only assumption). Appendix~\ref{app:rerun} lists every other pair count.

\newpage
Two replay findings held up under our robustness checks. First, the remaining excess loss is concentrated early
and on hard pairs (Figure~\ref{fig:excess}): 71\% of the label-removable excess (the part of the loss above $a(1-a)$
that more labels could remove; \invref, Proposition~4) sits at budgets 0--1, and
56\% sits on pairs with accuracy $\le 0.15$; whether its cause is the prior centre or the update strength is not
settled. Second, no widening of $\tau$ or of the
prior scale, and not the mixture prior, gives a gain whose interval excludes zero, either in the platform replay or in
a replay over 89 public benchmarks: the 92 benchmarks of \S\ref{sec:sim} minus the three original folds among them,
on which the shipped parameters were fitted.

% ---------------------------------------------------------------------------------------------------------------
\section{Simulator validity}\label{sec:sim}
A pre-registered check asked whether our local episodes look like the platform's, and it failed. Reading the
platform's per-pair feedback through the identity of \S\ref{sec:inversion} showed that platform pairs are harder and
more extreme than our folds: median pair accuracy 0.27 against 0.41 locally, 14.7\% of platform pairs at accuracy
$\le 0.05$, and none near 1.

The repair drew from the 92 benchmarks of the enlarged public pool that can produce platform-like episodes, and the
check passed: 0 of 11 platform runs outside their 90\% interval, a tail count of 6 within the allowed 1--7, and an
envelope test $p=0.544$. The pass survives dropping any single benchmark but not dropping all six largest
contributors, so we adopt it only in a narrowed scope: it validates the engine path without item groups, as
calibration rather than confirmation, and it predicts a repeat SD of 0.030 against 0.013 observed for the
hedged-centre model. An earlier
analysis used $p_0=0.5513$ instead of $0.448722$ and wrongly concluded that the platform was \emph{easier}.

% ---------------------------------------------------------------------------------------------------------------
\section{Headroom: floor and ceiling bounds}\label{sec:headroom}
\begin{figure}[H]
  \centering
  \begin{subfigure}[t]{0.49\linewidth}\centering
    \fig{fig_rho_value.pdf}
    \caption{Value of a label-free difficulty signal of quality \rs{}, computed with four named item fields as groups
    (not the shipped all-field grouping). 95\% bands; curve: $-71\rs^2$ fit.}
    \label{fig:rho}
  \end{subfigure}\hfill
  \begin{subfigure}[t]{0.49\linewidth}\centering
    \fig{fig_excess_map.pdf}
    \caption{Where the hedged-centre, item-group and final models lose on the platform beyond the pair-constant floor
    $a(1-a)$, by pair-accuracy band and budget (in-sample; slot counts per band in parentheses).}
    \label{fig:excess}
  \end{subfigure}
  \caption{(a) Headroom from item-level signal, measured locally. (b) Platform loss decomposition.
  \src{(a) \p{src/paiec/analysis/e6_headroom.py}; \p{experiments/runs/e6/e6.json}. (b)
  \p{experiments/pattern_mining/meta_2026-10-16/m10_excess_map.py}; \p{m10_excess_map.csv}}}
\end{figure}
\paragraph{What item knowledge is worth.} Revealing one true quantity at a time to the hedged-centre model shows that
item difficulty accounts for about half of the closable budget-0 gap. We had pre-registered an expectation of 60--80\%
for the benchmark offset. The value of a label-free difficulty signal is roughly quadratic in its quality over
$\rs \in [0.2, 1.0]$ ($-\dalc/\rs^2 = 58$--$74$; Figure~\ref{fig:rho}). A scorer at $\rs=0.3$ is worth
$5.7\ee{3}$. That is within the range our offline LLM rating reached (0.19--0.34), and above the 0.14--0.23
reported for unseen benchmarks~\cite{pacchiardi2024}. A runtime LLM item prior, which the rules permit, is therefore
the most promising untested direction; we chose a dependency-free, deterministic predictor as a design choice. A scorer must reach $\rs\approx0.5$ to be worth 0.018.

\newpage
Used
without shrinkage the same signal is harmful below $\rs\approx0.5$--$0.55$, so any external signal needs a learned weight.

\paragraph{Floors.} With the true offset, ability and item difficulty all supplied, budget-0 Brier is 0.1016. This is
an optimistic bound on one budget, not an ALC floor. A perfect item oracle reaches ALC 0.1387 locally, and nothing we
measured shows that the remaining loss is pure noise. On the platform, the \emph{pair-constant floor} $a(1-a)$, the
Brier score of knowing each pair's accuracy exactly but nothing about its items, is 0.151 on our saved draws (95\%
interval clustered by benchmark ID $[0.090, 0.205]$). Locally, the
same floor on the four folds is 0.1797 ($t$-interval $[0.135, 0.224]$), against 0.2088 for the final model.

\paragraph{Ceiling for pair-level methods.} The platform floor exceeds the mean ALC of 0.11 that our
pre-registration aimed for, but its benchmark-clustered interval includes it; so only on the observed draws does
0.11 require item-level knowledge. On those draws, the
largest gain any pair-level adjustment could give is $-0.0218$ ALC $[-0.0327, -0.0130]$, pooled in-sample
over three versions with identical budget-0 predictions. Locally,
the gap to the floor is $-0.0291$.

% ---------------------------------------------------------------------------------------------------------------
\section{Limitations}\label{sec:limits}
\begin{itemize}[leftmargin=*,nosep]
  \item \textbf{Evidence base.} Four held-out benchmarks: the item-group step's $t$-interval over fold means,
    $[-1.9, +0.2]\ee{3}$ (3 degrees of freedom), includes zero, and with one subject per benchmark cross-subject
    pooling is inactive. The simulator is valid only for the engine path without item groups, as calibration, and
    overstates the repeat SD. The oracle of \S\ref{sec:headroom} adds independent noise to true difficulty; a real
    scorer may be worth less.
  \item \textbf{Scope, registration and selection.} The out-of-sample benchmarks have no item features, so the item-group step
    is unconfirmed there. We did not test runtime LLM priors. 38 of the 60 logged experiments were registered after
    the fact; the offline LLM rating may have seen the public problems in training. We selected the final model (run 971666).
\end{itemize}

\begin{ack}
We thank the organisers for the competition, the data and answering our questions. Code, analyses and this report's
text were developed with the help of an LLM-based coding assistant, under the author's direction and review.
\end{ack}

% ---------------------------------------------------------------------------------------------------------------
\begin{thebibliography}{99}\small
\bibitem{aims2026} Stanford AIMS Lab. Predictive AI Evaluation Competition, NeurIPS 2026.
  \url{https://aims.stanford.edu/competition}.
\bibitem{truong2026} S.~Truong and S.~Koyejo. AI Measurement Science. Stanford AIMS Lab, 2026.
\bibitem{polo2024} F.~Maia Polo, L.~Weber, L.~Choshen, Y.~Sun, G.~Xu, and M.~Yurochkin. tinyBenchmarks: evaluating
  LLMs with fewer examples. In \emph{ICML}, 2024.
\bibitem{kipnis2025} A.~Kipnis, K.~Voudouris, L.~M.~Schulze Buschoff, and E.~Schulz. metabench: a sparse benchmark of
  reasoning and knowledge in large language models. In \emph{ICLR}, 2025.
\bibitem{vivek2024} R.~Vivek, K.~Ethayarajh, D.~Yang, and D.~Kiela. Anchor Points: benchmarking models with much fewer
  examples. In \emph{EACL}, pages 1576--1601, 2024.
\bibitem{truong2025} S.~Truong, Y.~Tu, P.~Liang, B.~Li, and S.~Koyejo. Reliable and efficient amortized model-based
  evaluation. In \emph{ICML}, 2025.
\bibitem{truong2026irsl} S.~Truong, Y.~Tu, R.~Schaeffer, and S.~Koyejo. Item response scaling laws: a measurement
  theory approach for efficient and generalizable neural scaling estimation. arXiv:2606.07616, 2026.
\bibitem{hofmann2025} V.~Hofmann, D.~Heineman, I.~Magnusson, K.~Lo, J.~Dodge, et~al. Fluid language model
  benchmarking. In \emph{COLM}, 2025.
\bibitem{shang2026} E.~Shang, W.~Tang, and Y.~He. LLM evaluation on unseen questions: contextual multidimensional IRT
  model. arXiv:2608.22295, 2026.
\bibitem{pacchiardi2024} L.~Pacchiardi, L.~G.~Cheke, and J.~Hern\'andez-Orallo. 100 instances is all you need:
  predicting the success of a new LLM on unseen data by testing on a few instances. KDD Workshop on Evaluation and
  Trustworthiness of Generative AI Models, 2024.
\bibitem{pacchiardi2025} L.~Pacchiardi, K.~Voudouris, B.~Slater, F.~Mart\'inez-Plumed, J.~Hern\'andez-Orallo,
  L.~Zhou, and W.~Schellaert. PredictaBoard: benchmarking LLM score predictability. In \emph{Findings of ACL}, 2025.
\bibitem{park2025} J.~Park, E.~Mendes, G.~Stanovsky, and A.~Ritter. Anticipatory evaluation of language models.
  arXiv:2509.20645, 2025.
\bibitem{ruan2024} Y.~Ruan, C.~J.~Maddison, and T.~Hashimoto. Observational scaling laws and the predictability of
  language model performance. In \emph{NeurIPS}, 2024.
\bibitem{efron1975} B.~Efron and C.~Morris. Data analysis using Stein's estimator and its generalizations.
  \emph{Journal of the American Statistical Association}, 70(350):311--319, 1975.
\bibitem{brown2008} L.~D.~Brown. In-season prediction of batting averages: a field test of empirical Bayes and Bayes
  methodologies. \emph{Annals of Applied Statistics}, 2(1):113--152, 2008.
\bibitem{gelman2008} A.~Gelman, A.~Jakulin, M.~G.~Pittau, and Y.-S.~Su. A weakly informative default prior
  distribution for logistic and other regression models. \emph{Annals of Applied Statistics}, 2(4):1360--1383, 2008.
\bibitem{brier1950} G.~W.~Brier. Verification of forecasts expressed in terms of probability. \emph{Monthly Weather
  Review}, 78(1):1--3, 1950.
\bibitem{gneiting2007} T.~Gneiting and A.~E.~Raftery. Strictly proper scoring rules, prediction, and estimation.
  \emph{Journal of the American Statistical Association}, 102(477):359--378, 2007.
\bibitem{murphy1973} A.~H.~Murphy. A new vector partition of the probability score. \emph{Journal of Applied
  Meteorology}, 12(4):595--600, 1973.
\bibitem{guo2017} C.~Guo, G.~Pleiss, Y.~Sun, and K.~Q.~Weinberger. On calibration of modern neural networks. In
  \emph{ICML}, PMLR 70:1321--1330, 2017.
\bibitem{kossen2021} J.~Kossen, S.~Farquhar, Y.~Gal, and T.~Rainforth. Active testing: sample-efficient model
  evaluation. In \emph{ICML}, 2021.
\bibitem{zrnic2024} T.~Zrnic and E.~J.~Cand\`es. Active statistical inference. In \emph{ICML}, 2024.
\bibitem{vanderlinden2000} W.~J.~van der Linden and C.~A.~W.~Glas. Capitalization on item calibration error in
  adaptive testing. \emph{Applied Measurement in Education}, 13(1):35--53, 2000.
\bibitem{bowyer2025} S.~Bowyer, L.~Aitchison, and D.~R.~Ivanova. Position: don't use the CLT in LLM evals with fewer
  than a few hundred datapoints. In \emph{ICML}, 2025.
\bibitem{miller2024} E.~Miller. Adding error bars to evals: a statistical approach to language model evaluations.
  arXiv:2411.00640, 2024.
\bibitem{wang2025} S.~Wang. Measuring all the noises of LLM evals. arXiv:2512.21326, 2025.
\end{thebibliography}

% ---------------------------------------------------------------------------------------------------------------
\appendix
\section{Reproducibility and rerun instructions}\label{app:rerun}
Everything runs on CPU. Our environment was Python 3.12.6 on Windows 11 on one laptop (AMD Ryzen 5 5500U, 6 cores,
5.8\,GB RAM), with packages pinned in
\texttt{requirements-lock.txt}. Code release (MIT licence): \url{https://github.com/Shashwat-srivastav/s2wish-paiec2026} (the release
commit is recorded in its history; every released file and its SHA-256 are listed in \texttt{RELEASE\_MANIFEST.txt}).
Seeds are fixed in each module.

\paragraph{Data.} The public training data are gated: request access to \p{aims-foundations/measurement-db} and
place the tables under \texttt{data/} at the revisions listed in Appendix~\ref{app:sources}. The simulator repair and
the 89-benchmark replay also need the 281-benchmark release. Our own formative feedback
(\texttt{experiments/platform\_submissions/}) is not part of the public release. The test of the re-implemented evaluator needs
the organisers' official baselines at the revision pinned in \texttt{external/PINS.txt}; the engine-equivalence tests need the gated data; the regression
tests of the feedback inversion need our feedback files. Tests whose inputs are absent skip rather than fail
(\texttt{REPRODUCE.md}, \S3).

\paragraph{Tests.} \texttt{python -m pytest -q} runs in about 7 minutes on CPU: 66 tests with the gated data and the
organisers' starting kit; 26 in a public clone (18 pass, 8 skip with reasons). The equivalence of the
re-implemented evaluator (\S\ref{sec:task}) is tested by \texttt{tests/test\_simulator.py}, which needs the starting kit. The suite also checks,
in \texttt{tests/test\_v1\_3\_archive.py}, that the final archive
has the quoted checksum, that its code files are byte-identical to the item-group model's, and that its predictions
match the research engine's all-field configuration on real episodes (maximum difference below $10^{-9}$).

\paragraph{Platform pair sets.} Several pair counts appear in the text; each comes from
\texttt{experiments/runs/review\_round1/pair\_counts.json} (singleton: the pair's benchmark appears once in its draw;
exact: every budget's success count is recovered within $2\varepsilon(p_0)$, uniquely and with nested prefixes).
\begin{center}\footnotesize
\begin{tabular}{@{}rp{0.52\linewidth}l@{}}
\toprule
size & pair set & used in\\
\midrule
8  & \emph{exact set}: replay-set pairs where the item-group and final models also replay exactly & \S\ref{sec:results}, Table~\ref{tab:platform}\\
14 & exact distinct singleton pairs among the 18 of row (b) & \invref; visual atlas\\
17 & the pairs of row (b) within a 0.005 tolerance & visual atlas (replay figure)\\
18 (a) & \emph{replay set}: exact distinct singleton pairs over our five saved formative runs & \S\ref{sec:inversion}, \S\ref{sec:results}, Table~\ref{tab:platform}\\
18 (b) & distinct singleton pairs of our later saved runs (14 exact) & this table\\
21 & exact singleton run-rows over our five saved formative runs & \invref\\
22 & distinct singleton pairs over our five saved formative runs & this table (denominator of row 18 (a))\\
26 & pair-slots of the three later models on the earlier draws (floor, ceiling) & \S\ref{sec:headroom}, Fig.~\ref{fig:excess}\\
27 & singleton run-rows over our five saved formative runs & \invref\\
34 & saved per-pair feedback slots (accuracy recovery check) & \S\ref{sec:inversion}\\
\bottomrule
\end{tabular}
\end{center}

\IfFileExists{appendix_repro_table.tex}{% Reproduction command table; mirrors REPRODUCE.md section 3 (maintained by hand).
% Runtimes: "check" measured 2026-10-19 in a clean-clone venv (Python 3.12.6, 3 BLAS threads); reruns = elapsed time
% recorded by the original run. Data: G = gated measurement-db under data/; F = our formative feedback (not in the
% public release); -- = public repository only. Column widths sum to 0.86\linewidth + 6\tabcolsep < \linewidth.
\begingroup
\footnotesize
\newcommand{\reprocmd}[1]{\texttt{#1}}
\newcommand{\reproab}{\allowbreak}
\renewcommand{\arraystretch}{1.05}
\begin{center}
\begin{tabular}{@{}>{\raggedright\arraybackslash}p{0.20\linewidth}>{\raggedright\arraybackslash}p{0.48\linewidth}%
>{\raggedright\arraybackslash}p{0.09\linewidth}>{\raggedright\arraybackslash}p{0.09\linewidth}@{}}
\toprule
item & command & runtime & data\\
\midrule
all checks, no data & \reprocmd{python scripts/\reproab reproduce\_all.py -{}-smoke} & 12\,s & --\\
all checks & \reprocmd{python scripts/\reproab reproduce\_all.py -{}-quick} & 82\,s & G, F\\
Table~\ref{tab:data} & \reprocmd{reproduce\_all.py -{}-only T1} & 2\,s & G\\
Table~\ref{tab:repr} & \reprocmd{reproduce\_all.py -{}-only T2}; rerun \reprocmd{python -m paiec.\reproab analysis.\reproab item\_repr} & ${<}1$\,s & --; rerun G\\
Table~\ref{tab:ladder} (R2, seeds 0--9) & \reprocmd{reproduce\_all.py -{}-only T3,\reproab T3live} & 67\,s & G\\
Table~\ref{tab:ladder}, official row & \reprocmd{python scripts/\reproab table3\_official\_baseline.py} & 30\,s & G\\
Table~\ref{tab:ladder}, R1 value & \reprocmd{reproduce\_all.py -{}-full -{}-jobs ladder\_r1} & not timed & G\\
Figures \ref{fig:teaser}--\ref{fig:forest} & \reprocmd{python docs/\reproab publication/\reproab neurips/\reproab make\_figures.py} & 3\,s & Fig.~\ref{fig:runs}: F\\
Table~\ref{tab:tried}, Fig.~\ref{fig:forest} reruns & \reprocmd{reproduce\_all.py -{}-full -{}-jobs exp0042,\reproab exp0046,\reproab exp0052,\reproab exp0053,\reproab exp0057} & 88\,min & G\\
Figure~\ref{fig:rho} rerun & \reprocmd{reproduce\_all.py -{}-full -{}-jobs e6} & 8\,min & G\\
Simulator check & \reprocmd{reproduce\_all.py -{}-only SIM}; rerun \reprocmd{python -m paiec.\reproab analysis.\reproab exp\_0039r\_\reproab simulator\_repair} & ${<}1$\,s & --; rerun G, F\\
Out-of-sample check (EXP-0058) & \reprocmd{reproduce\_all.py -{}-only EXP0058}; rerun \reprocmd{-{}-full -{}-jobs exp0058} & ${<}1$\,s; rerun 90\,s, 242\,MB & --; rerun G\\
Final archive rebuild &\reprocmd{python scripts/\reproab build\_v1\_3\_archive.py} (byte-identical) & ${<}1$\,s & --\\
Test suite (66 tests) & \reprocmd{python -m pytest -q} & 7\,min & partly G, F\\
\bottomrule
\end{tabular}
\end{center}
\endgroup
}%
  {\par\fillin{command-to-result table (appendix\_repro\_table.tex)}\par}

\paragraph{Archives.} SHA-256 checksums:\\[2pt]
{\footnotesize\begin{tabular}{@{}l@{}}
final model (selected; run 971666), \texttt{v1\_3\_named\_facets\_noguard.zip}\\
\quad\texttt{a4b0e841d4a344ff207bed2c855342212e95f7d550195b667446ea4e92685b86}\\
\end{tabular}}\par\medskip
The final archive's \texttt{model.py} and \texttt{labeling.py} are byte-identical to the item-group model's, and only
two parameters differ: \texttt{facets} \texttt{true}$\to$\texttt{"named"} and \texttt{guard} $1\to0$. The model tests
\texttt{facets} for truthiness, so \texttt{"named"} means \emph{all} fields, despite the archive's name; the effective
change is only that the key guard is off. (The research engine's literal \texttt{"named"} mode, used for
Figure~\ref{fig:rho}, is a different configuration.) \texttt{scripts/build\_v1\_3\_archive.py} rebuilds it byte for
byte (with zlib 1.3.x; with another zlib the archive members are still byte-identical). The supplementary
research record gives each registered experiment's question, hypothesis, settings, decision rule, result and commands.

\section{Training sources}\label{app:sources}
\begin{itemize}[leftmargin=*,nosep,before=\raggedright]
  \item \texttt{aims-foundations/measurement-db}, revision \texttt{bc8204d811} (CC BY-SA 4.0, gated): all models and
    parameters, including the final archive.
  \item \texttt{aims-foundations/measurement-db}, revision \texttt{54074453da} (the 281-benchmark release): the
    expanded-pool prior experiment and the simulator repair. No shipped parameter comes from it.
  \item \texttt{measurement-db-embed} vectors, revision \texttt{fb3a155cce} (licence not stated by the provider): the
    item-representation study (Table~\ref{tab:repr}) only.
  \item \texttt{sentence-transformers/all-MiniLM-L6-v2} (snapshot \texttt{1110a243fd}) and \texttt{all-mpnet-base-v2}
    (snapshot \texttt{e8c3b32edf}), both Apache 2.0: the distillation experiment only.
  \item TypeSafe Jev 1.13.0, a hosted service (licence not stated by the provider): offline LLM item ratings and
    decision checks. It is not used at scoring time.
  \item Formative platform feedback: used to compare versions and for the diagnostics of
    \S\ref{sec:inversion}--\S\ref{sec:headroom}.
\end{itemize}
The final archive uses no pretrained model and no network access at scoring time. Table~\ref{tab:repr} gives the
item-representation study cited in \S\ref{sec:task}.

\begin{table}[h]
  \caption{\rs{} between predicted and Rasch item difficulty (ridge regression, mean over the five benchmark folds).
  The LLM rating was computed offline. The rules allow scoring-time calls to approved inference APIs; our shipped
  predictor makes no network calls by design (no API dependency, deterministic, fully rerunnable). \src{\p{src/paiec/analysis/item_repr.py}; \p{experiments/runs/phase1_item_repr/item_repr_results.csv}}}
  \label{tab:repr}
  \centering\small
  \begin{tabular}{lcc}
    \toprule
    item representation & \rs{} within a benchmark & \rs{} on an unseen benchmark\\
    \midrule
    length and format                        & 0.22 & $-0.01$\\
    TF-IDF                                   & 0.36 & $-0.08$\\
    \texttt{measurement-db-embed} vectors    & ---  & $-0.08^{\dagger}$\\
    LLM difficulty rating (offline)          & 0.22 & \phantom{$-$}0.23\\
    all combined                             & 0.41 & \phantom{$-$}0.18\\
    \bottomrule
  \end{tabular}\\[2pt]
  {\footnotesize $^{\dagger}$Three-benchmark subset; the inner cross-validation chose a constant.}
\end{table}

\section{Names used in this report and internal identifiers}\label{app:names}
The research record, archive file names and the visual atlas use internal identifiers. This table maps them to the
plain names of the main text.
\begin{center}\scriptsize
\begin{tabular}{@{}p{0.40\linewidth}p{0.22\linewidth}l@{}}
\toprule
name in this report & internal identifier & archive or record\\
\midrule
official baseline (empirical mean) & L0 & organisers' starter kit\\
Beta-binomial shrinkage & v0 & \texttt{v0.zip}\\
first model (+ benchmark offset, item pooling) & v1 & \texttt{v1.zip}\\
hedged-centre model & v1\_1 & \texttt{v1\_1.zip}\\
item-group model & v1\_2 & \texttt{v1\_2.zip}\\
identical re-package of the item-group model & v1\_3 (guard on) & \texttt{v1\_3\_named\_facets.zip}\\
\textbf{final model} & v1\_3 & \texttt{v1\_3\_named\_facets\_noguard.zip}\\
metadata/IRT hybrid & v1\_5 & \texttt{v1\_5\_hybrid\_experimental.zip}\\
damped early updates & v1\_6 & \texttt{v1\_6\_damped\_facets.zip}\\
cap on early jumps & v1\_7 & \texttt{v1\_7\_adaptive\_facets.zip}\\
mixture prior & v1\_8, EXP-0057 & \texttt{v1\_8\_ext10.zip}\\
\midrule
platform-like draw / one subject per benchmark / all subjects & regimes R1 / R2 / R3 & \\
budget-$n$ Brier score & $B_n$ & \\
item groups / key guard & facets / guard (G1 key guard) & model parameters\\
false-merge probability $r$; group variance share $\lambda$ & $\rho$; $\pi$ & model parameters\\
\midrule
item scorer; label-free benchmark level; nested meta-learned priors & EXP-0041; EXP-0043; EXP-0045 & research record\\
lexical residual smoothing; gp-IRT reimplementation & EXP-0042; EXP-0046 & research record\\
final-model replay against the research engine & EXP-0044 & research record\\
prior factorial (stopped after pilot) & EXP-0049 & research record\\
re-tuned shrinkage strength; temperature scaling; stress test & EXP-0052; EXP-0053; EXP-0054 & research record\\
simulator validity check; its repair & EXP-0039; EXP-0039R & research record\\
pre-registered out-of-sample check (\S\ref{sec:results}) & EXP-0058 & \p{experiments/runs/EXP-0058/}\\
first final-model run; failed-but-scored item-group run & Codabench 967620; 955650 & platform feedback\\
value of an item-difficulty signal (Figure~\ref{fig:rho}) & E6 & \texttt{experiments/runs/e6}\\
\bottomrule
\end{tabular}
\end{center}

\begin{table}[h]
  \caption{Table~\ref{tab:ladder} with all six label budgets (Brier score; one subject per benchmark, mean over the
  four held-out benchmarks, 10 seeds, paired).
  \src{\p{experiments/runs/review_round1/ladder_intervals.json} and \p{table3_baseline_seeds0-9.json}}}
  \label{tab:ladder-full}
  \centering\small\setlength{\tabcolsep}{4pt}
  \begin{tabular}{@{}lccccccc@{}}
    \toprule
    model & $B_0$ & $B_1$ & $B_3$ & $B_7$ & $B_{15}$ & $B_{31}$ & ALC\\
    \midrule
    official baseline       & 0.250 & 0.360 & 0.243 & 0.209 & 0.194 & 0.189 & 0.2451\\
    Beta-binomial shrinkage & 0.250 & 0.231 & 0.214 & 0.201 & 0.192 & 0.188 & 0.2114\\
    first model             & 0.247 & 0.229 & 0.213 & 0.200 & 0.191 & 0.187 & 0.2098\\
    hedged-centre model     & 0.246 & 0.229 & 0.213 & 0.200 & 0.191 & 0.187 & 0.2097\\
    item-group model        & 0.246 & 0.229 & 0.212 & 0.200 & 0.190 & 0.184 & 0.2089\\
    final model             & 0.246 & 0.229 & 0.212 & 0.199 & 0.189 & 0.183 & 0.2088\\
    \bottomrule
  \end{tabular}
\end{table}

\begin{figure}[h]
  \centering
  \fig{fig_platform_runs.pdf}
  \caption{Logged formative runs of the development ladder and the rejected variants (one point each, fresh draws;
  bars are column means). \emph{Development ladder}:
  the models of Table~\ref{tab:ladder}; the open circle is a functionally identical re-package, so the item-group
  column has $n=5$ (Table~\ref{tab:platform} counts the item-group archive only, $n=4$). \emph{Rejected
  variants}: a metadata/IRT hybrid, damped early updates, a cap on early jumps, the mixture prior. Runs are unpaired: the final model alone ranges from 0.164 to 0.216.
  \src{\p{docs/publication/neurips/make_figures.py}; \p{experiments/platform_submissions/*.json}}}
  \label{fig:runs}
\end{figure}

\begin{figure}[h]
  \centering
  \fig{fig_programme_forest.pdf}
  \caption{The five pre-registered candidates of Table~\ref{tab:tried}, each paired against the final model with one
  subject per benchmark (\dalc{}\,\ee{3}, 95\% interval over pairs; positive is worse). The dashed line is the
  pre-registered adoption bar ($-1.0$). None qualifies.
  \src{\p{make_figures.py}; \texttt{experiments/runs/preregistration/exp\_00\{42,46,52,53,57\}\_*.json}}}
  \label{fig:forest}
\end{figure}

% Appendix D: exact inversion and replay of per-pair formative feedback (written separately; defines \label{app:inversion}).
\IfFileExists{appendix_inversion.tex}{% Exact inversion of per-pair formative feedback. Input by main.tex; defines \label{app:inversion}.
% Every number below is asserted by tests/test_feedback_inversion.py (python -m pytest -q tests/test_feedback_inversion.py).
% Code: src/paiec/feedback/ (scoring, inversion, krecovery, replay, bounds, io). Usage: INVERSION_README.md.
\section{Exact inversion of per-pair formative feedback}\label{app:inversion}

The platform's formative feedback reports, for every evaluated subject--benchmark pair of our own submission, the
Brier score $B_b$ and the calibration error $E_b$ at each label budget $b\in\mathcal{B}=\{0,1,3,7,15,31\}$, rounded to
six decimals. This appendix shows that, under conditions we state and check, these numbers determine (a)~the pair's
accuracy, (b)~on singleton pairs, the number of successes among the revealed labels, and hence (c)~the exact score of
any pair-level change to the model, together with floor and ceiling bounds. The implementation is the NumPy package
\texttt{paiec.feedback}; each number quoted here is asserted by \texttt{tests/test\_feedback\_inversion.py}.

\paragraph{Setting.}
A pair has a fixed evaluation pool of scored response rows with outcomes $y_i\in\{0,1\}$, the same at every budget;
its \emph{accuracy} is $a=\bar y$, the mean outcome over the scored rows. Labels are revealed from a separate
acquisition pool; at budget $b$ the model has seen the first $b$ of them, of which $k_b$ are successes. The scoring
rules are those published by the organisers: $B_b=\frac1n\sum_i(p_i-y_i)^2$, and $E_b=\sum_j\frac{n_j}{n}|\bar p_j-\bar
y_j|$ over ten fixed equal-width bins $[0,0.1),\dots,[0.9,1]$ weighted by response count; the pair's ALC is
$\sum_b w_bB_b$ with $w=(0.1,0.2,0.2,0.2,0.2,0.1)$. We assume:
\begin{enumerate}[label=A\arabic*.,leftmargin=2.2em,itemsep=0pt,topsep=2pt]
  \item \emph{Constant budget-0 prediction:} the model predicts one value $p_0$ on every scored row of the pair. We
    obtain $p_0$ by calling the archive's own \texttt{predict()} with no labels (and check it is constant across items
    and subjects); it is never read off the feedback.
  \item \emph{Exact scoring:} the reported values are the rules above applied to binary outcomes.
  \item \emph{Reported precision:} each reported value has absolute rounding error at most $\eta=5\times10^{-7}$.
  \item \emph{Count-only prediction (for $k$ only):} at budget $b$ the prediction on every scored row is $p(k_b,b)$,
    a function of the success count alone. For our hierarchical models this holds on \emph{singleton pairs} (the
    benchmark appears once in the draw, so no item pooling across subjects) whose items trigger no item-level term.
  \item \emph{Nested prefixes (for $k$ only):} with accept-all acquisition the platform reveals labels in a fixed
    order, so $0\le k_{b'}-k_b\le b'-b$ for $b<b'$, and a pair has the same $k$ path in every run.
\end{enumerate}

\paragraph{Lemma 1 (constant-prediction identities).}
If a pair's prediction is a constant $q$, then $B=a(1-2q)+q^2=a(1-a)+(q-a)^2$ and $E=|q-a|$, whatever the bin edges.
\textit{Proof.} Since $y_i^2=y_i$, $\frac1n\sum_i(q-y_i)^2=q^2-2q\bar y+\bar y$, which rearranges to both forms. All
rows fall in one bin, whose gap is $|q-\bar y|$ with weight 1.\hfill$\square$

\paragraph{Proposition 1 (accuracy).}
Under A1--A2 and $p_0\neq\frac12$,
\begin{equation}\label{eq:inv-a}
a=\frac{B_0-p_0^2}{1-2p_0},\qquad\text{and independently}\qquad a\in\{p_0-E_0,\;p_0+E_0\}.
\end{equation}
Under A3 the first expression has error at most $\eta/|1-2p_0|$, so it and the nearer ECE root disagree by at most
$\varepsilon(p_0)=\eta/|1-2p_0|+\eta$. \textit{Proof.} Solve Lemma~1 at $b=0$ for $a$; the error bound is the derivative
$1/(1-2p_0)$ times $\eta$, plus $\eta$ for $E_0$.\hfill$\square$

The bound shows where the method is fragile: the condition number $1/|1-2p_0|$ is about 10 at our $p_0=0.448722$,
so $\varepsilon=5.38\times10^{-6}$ there ($3.00\times10^{-6}$ at the first model's $p_0=0.400099$), and it diverges at $p_0=\frac12$.

\paragraph{Proposition 2 (identifiability and the reflection pitfall).}
Without $p_0$, the pair $(B_0,E_0)$ identifies only $a(1-a)=B_0-E_0^2$, i.e.\ the unordered set $\{a,1-a\}$; the
points $(a,a\pm E_0)$ and $(1-a,1-a\pm E_0)$ all reproduce $(B_0,E_0)$. In particular the reflection
$\rho:(a,p_0)\mapsto(1-a,1-p_0)$ leaves both reported values unchanged, so substituting $1-p_0$ for $p_0$ in
\eqref{eq:inv-a} returns $1-a$ \emph{and still passes the ECE cross-check}. At $p_0=\frac12$, $B_0=\frac14$ for every
$a$ and only $|a-\frac12|$ is identified. \textit{Proof.} By Lemma~1, $B_0-E_0^2=a(1-a)$, which fixes $a$ up to
$a\leftrightarrow1-a$, and $p_0=a\pm E_0$. Under $\rho$, $(1-a)a+((1-p_0)-(1-a))^2=a(1-a)+(a-p_0)^2$ and
$|(1-p_0)-(1-a)|=|a-p_0|$.\hfill$\square$

\emph{A cautionary example (our own error).} Until 2026-10-15 we used $p_0=0.5513$, which is $1-0.448722$ to four
decimals, instead of calling the archives. Every recovered accuracy was reflected and the ECE check passed on every
pair. The reflected reading put the median pair accuracy of our 34 saved pair-slots at 0.73, with none at or below
0.05. The correct reading gives 0.27, with 14.7\% at or below 0.05: platform pairs are harder than our local folds,
not easier. Quantities symmetric in $a$, such as Brier scores, ALC and the floor below, were unaffected. The lesson is
general: the orientation of the inversion must come from the model, never from the feedback.

\paragraph{Proposition 3 (success counts on singleton pairs).}
Under A1--A5, $E_b=|p(k_b,b)-a|$, so $k_b$ minimises the residual $r_b(k)=\bigl||p(k,b)-a|-E_b\bigr|$ over
$k=0,\dots,b$, where $p(k,b)$ is computed by calling the model with $k$ successes among $b$ synthetic labels. We accept
a recovery as \emph{exact} if $\min_k r_b(k)\le2\varepsilon(p_0)$ and \emph{unique} if every other $k$ has residual above
$2\varepsilon(p_0)$. Recovery is not unique when two counts give predictions at the same distance on either side of $a$, or
predictions closer together than the tolerance. A5 provides two independent checks: the path must be nested, and
it must be identical across runs. Lemma~1 provides a third: $B_b-a(1-a)-E_b^2=0$ whenever the prediction is constant.
\textit{Proof.} Lemma~1 at budget $b$ with $q=p(k_b,b)$, and the error bound of Proposition~1 applied to $a$ and
$E_b$.\hfill$\square$

\paragraph{Corollary (exact zero-upload replay).}
Let $q(k,b)$ be any \emph{pair-level} rule, i.e.\ one value per pair from the revealed count. Given $a$ and the
recovered path $(k_b)$, its score on the platform pair is exactly $B_b(q)=a(1-a)+(q(k_b,b)-a)^2$, with ALC
$\sum_bw_bB_b(q)$. No upload, rerun or fitting is involved. At budget 0 only $a$ is needed, so budget-0
counterfactuals are exact on every pair where the observed and the alternative budget-0 predictions are constant.

\paragraph{Proposition 4 (floor and ceiling).}
For any pair-level method, $B_b\ge a(1-a)$, with equality if and only if $q=a$. Hence the mean ALC over pairs is at
least the \emph{pair-constant floor} $F=\overline{a(1-a)}$, which the feedback identifies even without $p_0$ (as
$B_0-E_0^2$). With the \emph{label-removable excess} $X_b=\overline{B_b-a(1-a)}$ of the observed curves, no pair-level
change can improve ALC by more than $\sum_bw_bX_b=\overline{\mathrm{ALC}}-F$. \textit{Proof.} Lemma~1 and
$\sum_bw_b=1$.\hfill$\square$ Methods whose predictions vary across items within a pair, such as item pooling across
subjects, are not covered and can go below $F$.

\paragraph{Validation.}
\emph{Synthetic.} We simulated pairs, scored them with the published rules (which match our re-implemented evaluator to
$10^{-12}$), rounded to six decimals and inverted them. Over 5{,}000 draws, excluding those with $|1-2p_0|<0.02$, the
error of \eqref{eq:inv-a} never exceeded $\eta/|1-2p_0|$. With beta-binomial pair-level rules we generated 1{,}310
pairs (6{,}550 budget steps). The true $k_b$ was always within tolerance, and 6{,}549 steps had a unique minimiser,
which was always correct (the one non-unique step was a tie between two counts). Both self-replay and replay of an
alternative rule matched direct simulation within $2\varepsilon$. Item-varying predictions were flagged by the constancy
residual.
\emph{Our saved feedback} (Table~\ref{tab:inv-val}). The 34 saved pair-slots come from four runs of four archives.
The exact $k$ paths of 18 distinct singleton pairs (21 run-rows of 27 singleton rows; 14 pairs
without the first model's run) satisfied every check: residual below $5\times10^{-6}$, nested, identical in the 3
pairs seen in two runs, and a self-replay that reproduces every reported $B_b$ within $3.2\times10^{-6}$. The
remaining singleton rows had residuals of at least $5.1\times10^{-4}$, two orders of magnitude above the largest
exact residual, which is the expected footprint of item-level terms violating A4. We therefore do not use them.

\begin{table}[h]
\centering\footnotesize
\caption{Validation of the inversion on our own saved feedback. Bounds are from Proposition~1 with $\eta=5\times10^{-7}$.}
\label{tab:inv-val}
\begin{tabular}{@{}p{0.47\linewidth}p{0.22\linewidth}p{0.25\linewidth}@{}}
\toprule
quantity & value & check\\
\midrule
max $|a-a_{\mathrm{ECE}}|$, 34 saved pair-slots & $5.05\times10^{-6}$ & $\le\varepsilon=5.38\times10^{-6}$\\
$B_0$ reconstructed from $(a,p_0)$, 34 slots & error $2.8\times10^{-17}$ & Lemma~1\\
pair-constant floor $F$, 26 slots of the three later models (v1\_1--v1\_3) & 0.1511 & Prop.~4\\
label-removable excess $X_0,\dots,X_{31}$ (same 26) & 0.081, 0.037, 0.021, 0.007, 0.004, 0.000 & Prop.~4\\
ceiling for any pair-level change (same 26) & $-0.0218$ ALC & observed ALC 0.173\\
\bottomrule
\end{tabular}
\end{table}

\paragraph{Limitations.}
(i)~Every replay is in-sample. It scores a change on the draw that informed it, and platform draws differ (the
floor is 0.151 on the earlier draws). Replays describe a draw; they do not predict the summative score.
(ii)~Recovering $k$ is exact only on singleton pairs with count-only predictions; on our runs this leaves 18 pairs, so
pair-level conclusions rest on small samples. (iii)~Exactness is exactness up to the
reported precision. The bounds assume six-decimal rounding and degrade as $1/|1-2p_0|$, and a platform that rounds
more coarsely, or a model with $p_0$ near $\frac12$, would loosen or void them. (iv)~The method needs the submitted
model to be callable offline, to obtain $p_0$ and $p(k,b)$. It does not apply to item-varying models without
further assumptions.

\paragraph{Ethics and competition rules.}
The method reads only the formative feedback that the platform shows a team about its own submissions, together with
that team's own archives. It never accesses any feedback other than our own, nor the hidden data, and it never tries to determine
which real subject or benchmark lies behind an anonymous ID. Throughout, the IDs serve only as opaque grouping labels,
and we name none. Every use of feedback in this report is diagnostic.
}{}

% NeurIPS paper checklist (after references and appendices; does not count toward the page limit).
\clearpage
\IfFileExists{checklist.tex}{\section*{NeurIPS Paper Checklist}



\begin{enumerate}

\item {\bf Claims}
    \item[] Question: Do the main claims made in the abstract and introduction accurately reflect the paper's contributions and scope?
    \item[] Answer: \answerYes{}
    \item[] Justification: The abstract and introduction state the main finding and its scope: shrinkage carries most of the gain, the refinements are not distinguishable from zero across the four held-out benchmarks, and the remaining gap is item-level (\S\ref{sec:results}, \S\ref{sec:headroom}).    \item[] Guidelines:
    \begin{itemize}
        \item The answer \answerNA{} means that the abstract and introduction do not include the claims made in the paper.
        \item The abstract and/or introduction should clearly state the claims made, including the contributions made in the paper and important assumptions and limitations. A \answerNo{} or \answerNA{} answer to this question will not be perceived well by the reviewers. 
        \item The claims made should match theoretical and experimental results, and reflect how much the results can be expected to generalize to other settings. 
        \item It is fine to include aspirational goals as motivation as long as it is clear that these goals are not attained by the paper. 
    \end{itemize}

\item {\bf Limitations}
    \item[] Question: Does the paper discuss the limitations of the work performed by the authors?
    \item[] Answer: \answerYes{}
    \item[] Justification: Section~\ref{sec:limits} (Limitations) covers simulator validity, the four folds, platform noise, oracle assumptions, the scope of the out-of-sample check and post-hoc registration; the unreached 0.11 aspiration is discussed in \S\ref{sec:headroom}.
    \item[] Guidelines:
    \begin{itemize}
        \item The answer \answerNA{} means that the paper has no limitation while the answer \answerNo{} means that the paper has limitations, but those are not discussed in the paper. 
        \item The authors are encouraged to create a separate ``Limitations'' section in their paper.
        \item The paper should point out any strong assumptions and how robust the results are to violations of these assumptions (e.g., independence assumptions, noiseless settings, model well-specification, asymptotic approximations only holding locally). The authors should reflect on how these assumptions might be violated in practice and what the implications would be.
        \item The authors should reflect on the scope of the claims made, e.g., if the approach was only tested on a few datasets or with a few runs. In general, empirical results often depend on implicit assumptions, which should be articulated.
        \item The authors should reflect on the factors that influence the performance of the approach. For example, a facial recognition algorithm may perform poorly when image resolution is low or images are taken in low lighting. Or a speech-to-text system might not be used reliably to provide closed captions for online lectures because it fails to handle technical jargon.
        \item The authors should discuss the computational efficiency of the proposed algorithms and how they scale with dataset size.
        \item If applicable, the authors should discuss possible limitations of their approach to address problems of privacy and fairness.
        \item While the authors might fear that complete honesty about limitations might be used by reviewers as grounds for rejection, a worse outcome might be that reviewers discover limitations that aren't acknowledged in the paper. The authors should use their best judgment and recognize that individual actions in favor of transparency play an important role in developing norms that preserve the integrity of the community. Reviewers will be specifically instructed to not penalize honesty concerning limitations.
    \end{itemize}

\item {\bf Theory assumptions and proofs}
    \item[] Question: For each theoretical result, does the paper provide the full set of assumptions and a complete (and correct) proof?
    \item[] Answer: \answerYes{}
    \item[] Justification: The only formal results are the feedback-inversion identity and its corollaries; their assumptions (A1--A5) and proofs are in \invref.
    \item[] Guidelines:
    \begin{itemize}
        \item The answer \answerNA{} means that the paper does not include theoretical results. 
        \item All the theorems, formulas, and proofs in the paper should be numbered and cross-referenced.
        \item All assumptions should be clearly stated or referenced in the statement of any theorems.
        \item The proofs can either appear in the main paper or the supplemental material, but if they appear in the supplemental material, the authors are encouraged to provide a short proof sketch to provide intuition. 
        \item Inversely, any informal proof provided in the core of the paper should be complemented by formal proofs provided in appendix or supplemental material.
        \item Theorems and Lemmas that the proof relies upon should be properly referenced. 
    \end{itemize}

    \item {\bf Experimental result reproducibility}
    \item[] Question: Does the paper fully disclose all the information needed to reproduce the main experimental results of the paper to the extent that it affects the main claims and/or conclusions of the paper (regardless of whether the code and data are provided or not)?
    \item[] Answer: \answerYes{}
    \item[] Justification: Appendix~\ref{app:rerun} gives the environment, commands, seeds, data revisions and the checksum of every submitted archive, including the selected submission (the final model, Codabench run 971666); Appendix~\ref{app:sources} lists every training source; \S\ref{sec:method} gives all fitted parameters. The out-of-sample check EXP-0058 (\S\ref{sec:results}) was pre-registered with its benchmark list, arms, seeds and decision rule (\texttt{docs/preregistrations/EXP-0058\_oos\_confirmation\_preregistration.md}) and runs in about 90\,s.
    \item[] Guidelines:
    \begin{itemize}
        \item The answer \answerNA{} means that the paper does not include experiments.
        \item If the paper includes experiments, a \answerNo{} answer to this question will not be perceived well by the reviewers: Making the paper reproducible is important, regardless of whether the code and data are provided or not.
        \item If the contribution is a dataset and\slash or model, the authors should describe the steps taken to make their results reproducible or verifiable. 
        \item Depending on the contribution, reproducibility can be accomplished in various ways. For example, if the contribution is a novel architecture, describing the architecture fully might suffice, or if the contribution is a specific model and empirical evaluation, it may be necessary to either make it possible for others to replicate the model with the same dataset, or provide access to the model. In general. releasing code and data is often one good way to accomplish this, but reproducibility can also be provided via detailed instructions for how to replicate the results, access to a hosted model (e.g., in the case of a large language model), releasing of a model checkpoint, or other means that are appropriate to the research performed.
        \item While NeurIPS does not require releasing code, the conference does require all submissions to provide some reasonable avenue for reproducibility, which may depend on the nature of the contribution. For example
        \begin{enumerate}
            \item If the contribution is primarily a new algorithm, the paper should make it clear how to reproduce that algorithm.
            \item If the contribution is primarily a new model architecture, the paper should describe the architecture clearly and fully.
            \item If the contribution is a new model (e.g., a large language model), then there should either be a way to access this model for reproducing the results or a way to reproduce the model (e.g., with an open-source dataset or instructions for how to construct the dataset).
            \item We recognize that reproducibility may be tricky in some cases, in which case authors are welcome to describe the particular way they provide for reproducibility. In the case of closed-source models, it may be that access to the model is limited in some way (e.g., to registered users), but it should be possible for other researchers to have some path to reproducing or verifying the results.
        \end{enumerate}
    \end{itemize}


\item {\bf Open access to data and code}
    \item[] Question: Does the paper provide open access to the data and code, with sufficient instructions to faithfully reproduce the main experimental results, as described in supplemental material?
    \item[] Answer: \answerYes{}
    \item[] Justification: The code is public at \url{https://github.com/Shashwat-srivastav/s2wish-paiec2026}. The training data are gated by the organisers; Appendix~\ref{app:rerun} says how to request access and where to place them.
    \item[] Guidelines:
    \begin{itemize}
        \item The answer \answerNA{} means that paper does not include experiments requiring code.
        \item Please see the NeurIPS code and data submission guidelines (\url{https://neurips.cc/public/guides/CodeSubmissionPolicy}) for more details.
        \item While we encourage the release of code and data, we understand that this might not be possible, so \answerNo{} is an acceptable answer. Papers cannot be rejected simply for not including code, unless this is central to the contribution (e.g., for a new open-source benchmark).
        \item The instructions should contain the exact command and environment needed to run to reproduce the results. See the NeurIPS code and data submission guidelines (\url{https://neurips.cc/public/guides/CodeSubmissionPolicy}) for more details.
        \item The authors should provide instructions on data access and preparation, including how to access the raw data, preprocessed data, intermediate data, and generated data, etc.
        \item The authors should provide scripts to reproduce all experimental results for the new proposed method and baselines. If only a subset of experiments are reproducible, they should state which ones are omitted from the script and why.
        \item At submission time, to preserve anonymity, the authors should release anonymized versions (if applicable).
        \item Providing as much information as possible in supplemental material (appended to the paper) is recommended, but including URLs to data and code is permitted.
    \end{itemize}


\item {\bf Experimental setting/details}
    \item[] Question: Does the paper specify all the training and test details (e.g., data splits, hyperparameters, how they were chosen, type of optimizer) necessary to understand the results?
    \item[] Answer: \answerYes{}
    \item[] Justification: Folds, episode regimes and seeds are in \S\ref{sec:task}; priors, fitted values and integration in \S\ref{sec:method}; hyperparameters were fitted on the public benchmarks only (Appendix~\ref{app:sources}).
    \item[] Guidelines:
    \begin{itemize}
        \item The answer \answerNA{} means that the paper does not include experiments.
        \item The experimental setting should be presented in the core of the paper to a level of detail that is necessary to appreciate the results and make sense of them.
        \item The full details can be provided either with the code, in appendix, or as supplemental material.
    \end{itemize}

\item {\bf Experiment statistical significance}
    \item[] Question: Does the paper report error bars suitably and correctly defined or other appropriate information about the statistical significance of the experiments?
    \item[] Answer: \answerYes{}
    \item[] Justification: Comparisons are paired on identical episodes with 95\% bootstrap intervals over pairs; benchmark-level claims use $t$-intervals over the four fold means; platform runs are reported individually with their SDs (\S\ref{sec:task}, \S\ref{sec:results}). The pre-registered out-of-sample check EXP-0058 reports percentile benchmark-level bootstrap intervals over its 89 held-out benchmarks (10,000 resamples), with its primary and secondary outcomes fixed in advance (\S\ref{sec:results}). Bootstraps are percentile intervals; no normality is assumed except for the $t$-intervals over fold means and platform runs.
    \item[] Guidelines:
    \begin{itemize}
        \item The answer \answerNA{} means that the paper does not include experiments.
        \item The authors should answer \answerYes{} if the results are accompanied by error bars, confidence intervals, or statistical significance tests, at least for the experiments that support the main claims of the paper.
        \item The factors of variability that the error bars are capturing should be clearly stated (for example, train/test split, initialization, random drawing of some parameter, or overall run with given experimental conditions).
        \item The method for calculating the error bars should be explained (closed form formula, call to a library function, bootstrap, etc.)
        \item The assumptions made should be given (e.g., Normally distributed errors).
        \item It should be clear whether the error bar is the standard deviation or the standard error of the mean.
        \item It is OK to report 1-sigma error bars, but one should state it. The authors should preferably report a 2-sigma error bar than state that they have a 96\% CI, if the hypothesis of Normality of errors is not verified.
        \item For asymmetric distributions, the authors should be careful not to show in tables or figures symmetric error bars that would yield results that are out of range (e.g., negative error rates).
        \item If error bars are reported in tables or plots, the authors should explain in the text how they were calculated and reference the corresponding figures or tables in the text.
    \end{itemize}

\item {\bf Experiments compute resources}
    \item[] Question: For each experiment, does the paper provide sufficient information on the computer resources (type of compute workers, memory, time of execution) needed to reproduce the experiments?
    \item[] Answer: \answerYes{}
    \item[] Justification: Every experiment ran on one laptop CPU, an AMD Ryzen 5 5500U (6 cores, 12 threads) with 5.8\,GB of RAM visible to Windows~11, and no GPU (Appendix~\ref{app:rerun}). Measured wall times (\texttt{REPRODUCE.md}): the test suite takes 419\,s; checking every reported number from saved results takes 82\,s; rerunning the heavy experiments one at a time takes about 2.5\,h; the five experiments of Figure~\ref{fig:forest} took 1971, 353, 727, 1795 and 463\,s. The predictor takes 3.3--222\,ms per call (\S\ref{sec:method}). The total compute of the whole programme, including failed experiments, was not logged; from these runtimes we estimate it at tens of CPU-hours (an estimate, not a measurement).
    \item[] Guidelines:
    \begin{itemize}
        \item The answer \answerNA{} means that the paper does not include experiments.
        \item The paper should indicate the type of compute workers CPU or GPU, internal cluster, or cloud provider, including relevant memory and storage.
        \item The paper should provide the amount of compute required for each of the individual experimental runs as well as estimate the total compute. 
        \item The paper should disclose whether the full research project required more compute than the experiments reported in the paper (e.g., preliminary or failed experiments that didn't make it into the paper). 
    \end{itemize}
    
\item {\bf Code of ethics}
    \item[] Question: Does the research conducted in the paper conform, in every respect, with the NeurIPS Code of Ethics \url{https://neurips.cc/public/EthicsGuidelines}?
    \item[] Answer: \answerYes{}
    \item[] Justification: The work uses public, organiser-provided benchmark results; no human subjects or personal data are involved, and anonymous benchmark IDs were not de-anonymised (\invref).
    \item[] Guidelines:
    \begin{itemize}
        \item The answer \answerNA{} means that the authors have not reviewed the NeurIPS Code of Ethics.
        \item If the authors answer \answerNo, they should explain the special circumstances that require a deviation from the Code of Ethics.
        \item The authors should make sure to preserve anonymity (e.g., if there is a special consideration due to laws or regulations in their jurisdiction).
    \end{itemize}


\item {\bf Broader impacts}
    \item[] Question: Does the paper discuss both potential positive societal impacts and negative societal impacts of the work performed?
    \item[] Answer: \answerNo{}
    \item[] Justification: The paper has no separate impact section. Cheaper evaluation is the intended benefit; the main risk is using predicted scores in place of real evaluation, which could hide failures on a new benchmark. Our intervals and Limitations section quantify how uncertain such predictions are.
    \item[] Guidelines:
    \begin{itemize}
        \item The answer \answerNA{} means that there is no societal impact of the work performed.
        \item If the authors answer \answerNA{} or \answerNo, they should explain why their work has no societal impact or why the paper does not address societal impact.
        \item Examples of negative societal impacts include potential malicious or unintended uses (e.g., disinformation, generating fake profiles, surveillance), fairness considerations (e.g., deployment of technologies that could make decisions that unfairly impact specific groups), privacy considerations, and security considerations.
        \item The conference expects that many papers will be foundational research and not tied to particular applications, let alone deployments. However, if there is a direct path to any negative applications, the authors should point it out. For example, it is legitimate to point out that an improvement in the quality of generative models could be used to generate Deepfakes for disinformation. On the other hand, it is not needed to point out that a generic algorithm for optimizing neural networks could enable people to train models that generate Deepfakes faster.
        \item The authors should consider possible harms that could arise when the technology is being used as intended and functioning correctly, harms that could arise when the technology is being used as intended but gives incorrect results, and harms following from (intentional or unintentional) misuse of the technology.
        \item If there are negative societal impacts, the authors could also discuss possible mitigation strategies (e.g., gated release of models, providing defenses in addition to attacks, mechanisms for monitoring misuse, mechanisms to monitor how a system learns from feedback over time, improving the efficiency and accessibility of ML).
    \end{itemize}
    
\item {\bf Safeguards}
    \item[] Question: Does the paper describe safeguards that have been put in place for responsible release of data or models that have a high risk for misuse (e.g., pre-trained language models, image generators, or scraped datasets)?
    \item[] Answer: \answerNA{}
    \item[] Justification: We release a 5\,KB statistical predictor and analysis code; neither poses a high risk of misuse.
    \item[] Guidelines:
    \begin{itemize}
        \item The answer \answerNA{} means that the paper poses no such risks.
        \item Released models that have a high risk for misuse or dual-use should be released with necessary safeguards to allow for controlled use of the model, for example by requiring that users adhere to usage guidelines or restrictions to access the model or implementing safety filters. 
        \item Datasets that have been scraped from the Internet could pose safety risks. The authors should describe how they avoided releasing unsafe images.
        \item We recognize that providing effective safeguards is challenging, and many papers do not require this, but we encourage authors to take this into account and make a best faith effort.
    \end{itemize}

\item {\bf Licenses for existing assets}
    \item[] Question: Are the creators or original owners of assets (e.g., code, data, models), used in the paper, properly credited and are the license and terms of use explicitly mentioned and properly respected?
    \item[] Answer: \answerYes{}
    \item[] Justification: Every asset is credited with its revision and licence in Appendix~\ref{app:sources}. Licences as stated on the providers' Hugging Face cards (checked 2026-10-22): \texttt{aims-foundations/measurement-db}, CC BY-SA 4.0 (gated; used under its access terms and not redistributed); \texttt{sentence-transformers/all-MiniLM-L6-v2} and \texttt{all-mpnet-base-v2}, Apache 2.0. For \texttt{aims-foundations/measurement-db-embed} and the hosted TypeSafe Jev service, a licence is not stated by the provider.
    \item[] Guidelines:
    \begin{itemize}
        \item The answer \answerNA{} means that the paper does not use existing assets.
        \item The authors should cite the original paper that produced the code package or dataset.
        \item The authors should state which version of the asset is used and, if possible, include a URL.
        \item The name of the license (e.g., CC-BY 4.0) should be included for each asset.
        \item For scraped data from a particular source (e.g., website), the copyright and terms of service of that source should be provided.
        \item If assets are released, the license, copyright information, and terms of use in the package should be provided. For popular datasets, \url{paperswithcode.com/datasets} has curated licenses for some datasets. Their licensing guide can help determine the license of a dataset.
        \item For existing datasets that are re-packaged, both the original license and the license of the derived asset (if it has changed) should be provided.
        \item If this information is not available online, the authors are encouraged to reach out to the asset's creators.
    \end{itemize}

\item {\bf New assets}
    \item[] Question: Are new assets introduced in the paper well documented and is the documentation provided alongside the assets?
    \item[] Answer: \answerYes{}
    \item[] Justification: The released code is documented in its README and release manifest, and the feedback-inversion package has its own documentation and tests. Licence: the code release is under the MIT License (\texttt{LICENSE} in the repository).
    \item[] Guidelines:
    \begin{itemize}
        \item The answer \answerNA{} means that the paper does not release new assets.
        \item Researchers should communicate the details of the dataset\slash code\slash model as part of their submissions via structured templates. This includes details about training, license, limitations, etc. 
        \item The paper should discuss whether and how consent was obtained from people whose asset is used.
        \item At submission time, remember to anonymize your assets (if applicable). You can either create an anonymized URL or include an anonymized zip file.
    \end{itemize}

\item {\bf Crowdsourcing and research with human subjects}
    \item[] Question: For crowdsourcing experiments and research with human subjects, does the paper include the full text of instructions given to participants and screenshots, if applicable, as well as details about compensation (if any)? 
    \item[] Answer: \answerNA{}
    \item[] Justification: No crowdsourcing or human-subject research was involved.
    \item[] Guidelines:
    \begin{itemize}
        \item The answer \answerNA{} means that the paper does not involve crowdsourcing nor research with human subjects.
        \item Including this information in the supplemental material is fine, but if the main contribution of the paper involves human subjects, then as much detail as possible should be included in the main paper. 
        \item According to the NeurIPS Code of Ethics, workers involved in data collection, curation, or other labor should be paid at least the minimum wage in the country of the data collector. 
    \end{itemize}

\item {\bf Institutional review board (IRB) approvals or equivalent for research with human subjects}
    \item[] Question: Does the paper describe potential risks incurred by study participants, whether such risks were disclosed to the subjects, and whether Institutional Review Board (IRB) approvals (or an equivalent approval/review based on the requirements of your country or institution) were obtained?
    \item[] Answer: \answerNA{}
    \item[] Justification: No human subjects were involved.
    \item[] Guidelines:
    \begin{itemize}
        \item The answer \answerNA{} means that the paper does not involve crowdsourcing nor research with human subjects.
        \item Depending on the country in which research is conducted, IRB approval (or equivalent) may be required for any human subjects research. If you obtained IRB approval, you should clearly state this in the paper. 
        \item We recognize that the procedures for this may vary significantly between institutions and locations, and we expect authors to adhere to the NeurIPS Code of Ethics and the guidelines for their institution. 
        \item For initial submissions, do not include any information that would break anonymity (if applicable), such as the institution conducting the review.
    \end{itemize}

\item {\bf Declaration of LLM usage}
    \item[] Question: Does the paper describe the usage of LLMs if it is an important, original, or non-standard component of the core methods in this research? Note that if the LLM is used only for writing, editing, or formatting purposes and does \emph{not} impact the core methodology, scientific rigor, or originality of the research, declaration is not required.
    %this research? 
    \item[] Answer: \answerYes{}
    \item[] Justification: An LLM (TypeSafe Jev 1.13.0) produced offline item-difficulty ratings and decision checks used in our experiments (Table~\ref{tab:repr}, \S\ref{sec:headroom}, Appendix~\ref{app:sources}); it is not used at scoring time, by design (the rules permit approved inference APIs at scoring time; our predictor makes no network calls). An LLM-based coding assistant helped develop code and text (Acknowledgements).
    \item[] Guidelines:
    \begin{itemize}
        \item The answer \answerNA{} means that the core method development in this research does not involve LLMs as any important, original, or non-standard components.
        \item Please refer to our LLM policy in the NeurIPS handbook for what should or should not be described.
    \end{itemize}

\end{enumerate}}{}

\ifreference
\clearpage
% Generated by make_atlas.py (build.ps1 regenerates it before every compile). Reference build only.
\IfFileExists{atlas.tex}{\input{atlas}}{}
\fi

\end{document}
